\documentclass[../../main]{subfiles}

\begin{document}

\section{Propriedades das Relações}
\label{sec:matematica:relacoes:propriedades}

\begin{defn}[Relação Reflexiva]

	Seja \(R\subseteq A\times A\).

	Dizemos que \(R\) é reflexiva quando

	\[
		(\forall a\in A)
		(aRa).
	\]

\end{defn}

\begin{defn}[Relação Irreflexiva]

	Seja \(R\subseteq A\times A\).

	Dizemos que \(R\) é irreflexiva quando

	\[
		(\forall a\in A)
		\neg(aRa).
	\]

\end{defn}

\begin{defn}[Relação Simétrica]

	Seja \(R\subseteq A\times A\).

	Dizemos que \(R\) é simétrica quando

	\[
		(\forall a,b\in A)
		(aRb \Rightarrow bRa).
	\]

\end{defn}

\begin{defn}[Relação Antissimétrica]

	Seja \(R\subseteq A\times A\).

	Dizemos que \(R\) é antissimétrica quando

	\[
		(\forall a,b\in A)
		\bigl(
		(aRb)
		\land
		(bRa)
		\Rightarrow
		a=b
		\bigr).
	\]

\end{defn}

\begin{defn}[Relação Assimétrica]

	Seja \(R\subseteq A\times A\).

	Dizemos que \(R\) é assimétrica quando

	\[
		(\forall a,b\in A)
		(aRb \Rightarrow \neg(bRa)).
	\]

\end{defn}

\begin{defn}[Relação Transitiva]

	Seja \(R\subseteq A\times A\).

	Dizemos que \(R\) é transitiva quando

	\[
		(\forall a,b,c\in A)
		\bigl(
		(aRb)
		\land
		(bRc)
		\Rightarrow
		aRc
		\bigr).
	\]

\end{defn}

\begin{defn}[Relação Total]
	\label{def:matematica:relacoes:propriedades:total}

	Seja \(R\subseteq A\times A\).

	Dizemos que \(R\) é total quando

	\[
		(\forall a,b\in A)
		(a\neq b
		\Rightarrow
		(aRb \lor bRa)).
	\]

\end{defn}

\begin{prop}

	Toda relação assimétrica é irreflexiva.

\end{prop}

\begin{proof}

	Suponha que exista \(a\in A\) tal que \(aRa\).

	Pela assimetria,

	\[
		aRa
		\Rightarrow
		\neg(aRa),
	\]

	contradição.

\end{proof}

\end{document}
